\chapter{RQ1 --- The Minimal Sufficient Spatial Address}
\label{ch:rq1}

\emph{A type-conditional, image-recoverable key.}

\section{The question: what uniquely identifies an element among its look-alikes?}

The grounding task ends by selecting one IFC element from a retrieved candidate pool in which the
ground truth is always present (median 76, in-pool 100\%). The difficulty is not recall but
\emph{discrimination among visually identical siblings}: a curtain-wall fa\c{c}ade is a row of
near-identical windows; a floor is a run of indistinguishable doors. We formalize this as the
\textbf{confusable set}. For a target element $e$, let

\begin{quote}
$C(e)$ = the set of pool elements sharing $e$'s \textbf{coarse fingerprint} --- same building storey
and same IFC class --- i.e.\ the elements a coarse ontological filter cannot tell apart.
\end{quote}

RQ1 asks for the \textbf{minimal sufficient spatial address}: the smallest descriptor set that
uniquely identifies $e$ within $C(e)$, subject to two deployment constraints that make the address
\emph{usable} rather than merely \emph{distinguishing}. The address must be \textbf{IFC-computable}
(derivable from the BIM model with no human labelling, so it exists on day one of a cold-start
deployment) and \textbf{image-recoverable} (estimable from a site photograph and a plan, so a real
extractor can in principle recover it --- the constraint RQ2 then tests). A descriptor that
distinguishes but cannot be computed, or computes but cannot be seen, does not answer RQ1.

\section{The coarse floor saturates: ontology alone cannot separate siblings}

The natural first descriptor is the ontology itself --- storey and IFC class. It is necessary but
nowhere near sufficient. Restricting the pool to the target's storey and class leaves a \textbf{median
confusable set of 46} (mean 112) and an oracle Top-1 of only \textbf{4.9\%} --- \emph{below} the
realized G8 reranker (6.7\%), and essentially equal to it on Top-10 (oracle 31.5\% vs realized
30.0\%). The coarse fingerprint is \textbf{saturated}: supplying it at an oracle level buys nothing
the learned system does not already have, because it cannot separate elements that are, by
construction, of the same class on the same floor. Ontology is the \emph{floor}, not the
discriminator.

Adding the richest available \emph{attribute} --- the Revit family/type string \texttt{object\_type}
--- shrinks the confusable set to a \textbf{median of 13} (a $3.8\times$ reduction) and lifts oracle
Top-10 to 76\%, but it plateaus on Top-1: only \textbf{2 of 60} targets are uniquely identified by
attributes alone. Attributes narrow; they do not single out. What remains after attributes is the
irreducibly \emph{relational} residue --- which of the thirteen identical-looking siblings is
\emph{this} one --- and that residue is where the address must live.

\section{The headline finding: the address is type-conditional}

The central result of this chapter is that \textbf{the minimal sufficient address is not uniform
across element classes --- it is conditional on the element's type}, because what makes two elements
distinguishable differs by what kind of thing they are. Two descriptors, each the address for its
class, close the gap:

\begin{itemize}
  \item \textbf{Fillers (windows, doors)} are identified by their \textbf{position-slot} $(i, M)$:
  the target is the $i$-th of $M$ openings ordered along its host wall. A window is pinned not by its
  own appearance but by its ordinal place in a row.
  \item \textbf{Walls} are identified by a \textbf{connectivity fingerprint} ---
  (\texttt{connection\_degree}, \texttt{hosted\_opening\_count}, \texttt{length\_band},
  \texttt{is\_external}): how many walls it joins, how many openings it hosts, its length band, and
  whether it is on the building envelope. Within the same-storey walls, this collapses a median
  confusable set of $110 \to 2$ (uniquely identifying 10 of 22 walls, where \texttt{object\_type}
  identifies 0).
\end{itemize}

Supplied at an oracle level, this type-conditional address takes the \emph{real retrieved pool} from
a coarse Top-1 of 4.9\% to \textbf{78.5\%}, and Top-10 from 31.5\% to \textbf{98.1\%} (MRR $0.137 \to
0.854$), at zero recall cost --- a more than fifteen-fold lift on Top-1 over the coarse floor. The
gain decomposes exactly along the type-conditional split: \textbf{fillers reach Top-1 91.0\%}
(position-slot), \textbf{walls 64.2\%} (fingerprint). No single uniform descriptor achieves this; the
position-slot is meaningless for a wall, and the connectivity fingerprint is degenerate for a window.
The answer to ``is the address element-class-conditional?'' is an empirical \emph{yes}, and the
conditioning is the contribution.

\section{Why depth-1 suffices: the address is a node, not a path}

A relational address invites the question of \emph{how relational} --- how many hops of graph context
the representation needs. The measurement is decisive: the oracle confusable set shrinks from a median
of 13 (attributes) to \textbf{8.2 at one hop}, then to \textbf{8.1} and \textbf{8.1} at two and three
hops. Almost all realizable discrimination is recovered at \textbf{depth 1}, and deeper context adds
essentially nothing, because the per-hop reliability of recovering a \emph{named} multi-hop relation
from an image collapses (the same evidence that motivates extracting at depth $\le 1$ in RQ3). The
address therefore has two equivalent forms --- a \textbf{depth-1 sub-graph} (the target node, its
directly-related neighbours, and the edge types) and the \textbf{distilled \texttt{node.val}
attributes} hung on the target --- and the two carry the same information. We match on the node
attributes and derive them from the sub-graph; no deep graph and no learned graph fusion is required.
The relations needed are few: \texttt{FILLS} and \texttt{NEXT\_TO} for the filler address,
\texttt{CONNECTS\_TO} and the reverse of \texttt{FILLS} for the wall address, with \texttt{ADJACENT\_TO}
as generic proximity --- sufficient for the two solved classes.

\section{The address satisfies its own two constraints}

The two deployment constraints are not assumed; they are checked.

\noindent\textbf{IFC-computable.} Every address field is reconstructed directly from the raw IFC
model (\texttt{AdvancedProject.ifc}) via \texttt{ifcopenshell} --- the position-slot from the
\texttt{IfcRelFillsElement} / \texttt{IfcRelVoidsElement} chains, the wall fingerprint from
\texttt{IfcRelConnectsPathElements} and the hosted openings \citep{buildingsmart2024ifc4x3}. The
reconstruction is validated against the dataset's own skeleton: the recovered wall
\texttt{connection\_degree} matches the ground-truth skeleton in \textbf{14 of 14} checked cases. The
address exists for every element with no human annotation --- the cold-start property the system
claims.

\noindent\textbf{Image-recoverable --- and not automatically so.} Image-recoverability is a
\emph{design constraint on the address}, not a property we may assume, and the two solved classes fall
on opposite sides of it. The filler position-slot is recoverable: the plan's openings render as
coloured marks, and the slot is read by ordering them along the host wall --- provided the ordinal is
numbered by an \textbf{image-coordinate} convention rather than the wall's arbitrary IFC local axis
(the representational prerequisite quantified in RQ2). The wall fingerprint, by contrast, is largely
\emph{not} recoverable: its load-bearing fields --- \texttt{connection\_degree} and
\texttt{length\_band} --- depend on where the modeller split a run into IfcWall \emph{instances}, and
collinear adjacent instances render as one continuous wall poch\'e, so the instance boundaries those
fields require are simply not in the image (measured: length-band recovered on 5/17 walls, realized
Top-1 $\approx$ floor; M2b). Only the wall's hosted-opening count, which renders, is recoverable, and
it is too weak alone. A descriptor keyed to an invisible modelling choice is thus excluded from
realization by construction --- which is exactly why the realized system is built on the one
image-recoverable address (the filler slot), the wall remaining an oracle-level result.
Image-recoverability is therefore not an afterthought but a property the address must be \emph{designed}
to satisfy.

\section{The honest boundary of the claim}

Four caveats bound the representation result. \emph{Oracle level:} every number here assumes perfect
extraction ($r = 1$); it establishes the \emph{ceiling}, and what fraction is \emph{realizable} under
a real detector is the separate question RQ2 answers (fillers: oracle 91 $\to$ realized 67.6).
\emph{Class coverage:} the address is closed for fillers and walls --- the bulk of the targets --- but
\textbf{3 ``other''-class} targets and room/space addressing remain open; notably the model exposes no
\texttt{IfcRelSpaceBoundary} (rooms are unnamed in this synthetic export), so space-relative addressing
is future work, not a claim of universal completeness. \emph{Single project:} the ceiling is measured
on one synthetic AP model; the \emph{form} of the address (type-conditional, topology-derived) is
general, but the specific descriptor reliabilities are project-specific. \emph{Descriptor sweep:} we
report the position-slot and the connectivity fingerprint; a fuller sweep of candidate descriptors
(junction type, generated cell, landmark-relative) with per-descriptor extractability scoring is left
to future work.

\section{Conclusion}

RQ1 asked what minimally and sufficiently identifies an IFC element among its visual look-alikes,
under the constraints that the descriptor be model-computable and evidence-recoverable. The answer is
a \textbf{type-conditional spatial address}: a coarse ontological prefix (storey + class) that is
necessary but saturated, completed by a class-specific topological body --- the position-slot for
fillers, the connectivity fingerprint for walls. Supplied at an oracle level it reorders the candidate
pool from Top-1 4.9\% to \textbf{78.5\%} at no recall cost, with the gain partitioned cleanly by
element type; the discrimination saturates at one graph hop, so the address compiles into a node
rather than a path; and it is verified to be both IFC-computable (14/14 against the skeleton) and
image-recoverable (under the right ordinal convention). The representation is the thesis's primary
contribution --- \emph{what} the minimal sufficient address is, and \emph{that} it is class-conditional
--- and it stands independently of how reliably any particular detector recovers it.
